\paragraph{富集空间对照能识别与不能识别的因素。}
Legendre35 和 Legendre36 控制同样包含式~\eqref{eq:layer-features}中的六个解析层函数。因此，本实验并没有证明“无富集的 Legendre 成功而无富集的三角空间失败”。同一分辨率下两族使用同一套预设径向求积；选择 Legendre 试验基不会自动把它替换为该基族专有的 Gauss--Lobatto 网格。多项式插值节点的端点聚束因而不是本对照中实测改变的解释变量。无富集且匹配保留秩的多项式控制仍需单独实验。

层函数按 $e^{-\gamma(L-s)}$ 衰减，集中在\emph{外边界} $s=L$，而不是圆柱内壁 $s=0$。此外，当前远场齐次算子采用 $\chi=16z^2(1-z)$ 和式~\eqref{eq:boundary-transform}中的修正；其无滑移条件并不强迫内壁曲率为零。以撇号表示 $z$ 导数，有 $(\mathcal Bq)''(0)=32q(0)-[(\chi q)'']_{z=1}$，通常非零。因此，仅凭删除这些层函数不能建立“内壁曲率自由度丢失导致分离角改变”的因果链。末态大误差表明当前结果对试验空间和保留维数共同敏感；固定中点设计与五个抽样种子的最坏误差相当或稍好，也不是随机化永远无益的定理。
